% Глава 3. Материал: docs/requirements/quality-scenarios.md, docs/research/metrics-gqm.md,
% docs/research/preregistration.md, docs/data/annotation-guide.md, docs/data/reports.md,
% docs/architecture/design-iterations.md, docs/risks.md. Тема и RQ — ADR-0052, гипотезы — preregistration.md v0.4.
\chapter{Методология}\label{ch:method}

Глава объясняет, как работа создаёт и оценивает свои артефакты. Раздел~\ref{sec:method-dsr} вводит
исследовательскую парадигму. Разделы~\ref{sec:method-req}--\ref{sec:method-gqm} описывают, как получены
требования, архитектура и метрики. Разделы~\ref{sec:method-hyp}--\ref{sec:method-s4} представляют четыре
исследования, а разделы~\ref{sec:method-prov}--\ref{sec:method-validity} --- предрегистрацию, этику и
угрозы валидности. Раздел~\ref{sec:method-harness} описывает, как обвязка разработки защищает и код, и
эксперименты.

\section{Исследовательская парадигма}\label{sec:method-dsr}

Работа выполнена в парадигме design science research: знание получают, создавая артефакт, решающий
реальную проблему, и строго его оценивая \cite{SRC-100}. Виринга разделяет внутри такого проекта два вида
работы: \emph{проектные задачи}, которые спрашивают, как изменить мир («создать тренажёр, который…»), и
\emph{вопросы знания}, которые спрашивают, что в нём верно («как часто модель отклоняет верное
чтение?», «насколько точно модель оценивает усвоение?») \cite{SRC-101}. Исследовательские вопросы
работы --- вопросы знания; задачи З1--З5 --- проектные задачи, которые делают их отвечаемыми. Таблица~\ref{tab:hevner} показывает, как работа
выполняет семь принципов Хевнера.

\begin{table}[ht]
\centering\small
\caption{Принципы design science Хевнера \cite{SRC-100} в этой работе.}
\label{tab:hevner}
\begin{tabularx}{\linewidth}{@{}P{3.8cm}L@{}}
\toprule
Принцип & В этой работе \\
\midrule
Проектирование как артефакт & Артефакт A --- набор данных и инструмент сбора (тренажёр по книге); артефакт B
--- модель распознавания, которая учится у учителей (B1, B2); артефакт C --- модель усвоения навыков книги
(K1, K2). \\
Значимость проблемы & Начинающие занимаются одни; время учителя не масштабируется; нет инструмента,
проверяющего чтение букварного уровня на языке ученика. \\
Оценка проекта & Надёжность разметки и согласие модели с учителями (RQ-1), контролируемый эксперимент на
замороженном тесте (RQ-2), проверка предсказания и согласия с тестом учителя (RQ-3), лонгитюдное наблюдение
(RQ-4). \\
Исследовательский вклад & Набор данных; адаптация модели на неизученной аудитории; оценка усвоения по речи с
учётом ошибок распознавателя; кривые обучения учеников. \\
Строгость & Предрегистрация; разбиение по ученикам и слепой тест; стандартные метрики MDD; статистики согласия с доверительными интервалами;
протокол кейс-стади \cite{SRC-102}; эмпирические стандарты как чек-лист \cite{SRC-103}. \\
Проектирование как поиск & 50 записанных решений с отклонёнными вариантами; семь итераций
проектирования. \\
Сообщение результатов & Диссертация, карточка набора данных, дообученные веса, открытая архитектура,
план докторантуры. \\
\bottomrule
\end{tabularx}
\end{table}

\section{Требования}\label{sec:method-req}

\subsection{Аудитория как главное ограничение}

Одно ограничение было сформулировано первым и проверялось при каждом следующем решении: ученики ---
русскоязычные взрослые, которые учатся читать по-арабски с нуля. Из него следуют семь правил: объяснять
каждый знак по-русски до того, как он понадобится; идти в порядке книги маленькими шагами; давать
услышать до чтения; объяснять ошибки через знакомые звуки и показывать их в тексте; закреплять короткими
вопросами; никогда не наказывать; относиться к тексту Корана с уважением --- точный текст Мединского
мусхафа, начало с \emph{истиазы} и \emph{басмалы}.

\subsection{Функциональные требования и пользовательские истории}

Требования собраны из просьб автора, прототипа, книги и литературы и записаны с постоянными
идентификаторами и приоритетами MoSCoW. Текущая спецификация содержит 47 функциональных требований:
программа по книге (FR-01, FR-16, FR-36--FR-41), оценка и объяснения (FR-03--FR-05, FR-20--FR-22, FR-43),
опыт ученика (FR-02, FR-07, FR-31--FR-35), исследовательские данные (FR-06, FR-11--FR-14, FR-17, FR-24,
FR-26--FR-29), приватность (FR-08--FR-10, FR-44) и поздние дополнения интерфейса
(FR-45--FR-47). Пользовательские истории описывают те же требования
глазами ученика, учителя и исследователя. Матрица трассировки связывает каждое требование с кодом и
тестами; автоматический отчёт показывает, какие требования покрыты проходящими тестами (59 из 60 на момент
написания).

\subsection{Сценарии атрибутов качества}

Нефункциональные требования переписаны как сценарии атрибутов качества из шести частей --- источник,
стимул, окружение, артефакт, реакция, мера реакции --- по Бассу, Клементсу и Казману \cite{SRC-047}.
Приоритеты заданы деревом полезности по важности и техническому риску. Сценарии перечислены в
таблице~\ref{tab:qas}.

\begin{table}[ht]
\centering\small
\caption{Сценарии атрибутов качества (QAS-01--QAS-11).}
\label{tab:qas}
\begin{tabularx}{\linewidth}{@{}lP{3.8cm}L@{}}
\toprule
ID & Атрибут & Сценарий и мера \\
\midrule
QAS-01 & Точность & Вердикт совпадает с учителем: precision, recall, F1; $\kappa$ «бот --- учитель» против
«учитель --- учитель». \\
QAS-02 & Точность & Верное эталонное чтение засчитывается: доля «не засчитано» на эталонах
$\le$5\,\% (утвердить). \\
QAS-03 & Производительность & Ответ на голосовое: 90\,\% попыток обработаны за $\le$10\,с. \\
QAS-04 & Надёжность & Сбой модели не ломает урок; 100\,\% сбоев записаны и получают ответ. \\
QAS-05 & Надёжность & Место в книге переживает перезапуск: ни одной потерянной позиции. \\
QAS-06 & Устойчивость ввода & Слишком короткие, беззвучные и слишком длинные записи не оцениваются. \\
QAS-07 & Удобство & Не менее четырёх из пяти новичков получают засчитанную строку в первой сессии без
помощи (утвердить). \\
QAS-08 & Приватность & Без согласия аудио нет; после удаления не остаётся ни строк, ни файлов. \\
QAS-09 & Воспроизводимость & Любую попытку можно пересчитать: версии, пороги и хеш сохранены. \\
QAS-10 & Модифицируемость & Новый распознаватель добавляется без правки домена. \\
QAS-11 & Целостность процесса & Агент не может ослабить тесты в фазе реализации. \\
\bottomrule
\end{tabularx}
\end{table}

\section{Метод проектирования архитектуры}\label{sec:method-add}

Архитектура получена облегчённым вариантом Attribute-Driven Design \cite{SRC-047}: каждая итерация берёт
одну цель и несколько драйверов (сценарии, требования, ограничения), записывает решения как записи
архитектурных решений (ADR) \cite{SRC-010} и завершается тестом, подтверждающим решение. Для одного автора
полные документы итераций были бы бюрократией, поэтому итерация --- одна строка таблицы, а подробности,
включая отклонённые варианты, лежат в ADR. Выполнено семь итераций: заменяемые модель, база данных и
мессенджер; исследовательские данные с отдельно хранимыми персональными данными; оценка, которой можно
доверять и которую можно оспорить; программа строго по книге; интерфейс, понятный без учителя;
разработка с агентом, которую нельзя «подогнать»; диагностика без утечки персональных данных. Результаты
описаны в главе~\ref{ch:impl}.

\section{Система измерений}\label{sec:method-gqm}

Метрики выведены методом «цель --- вопрос --- метрика» \cite{SRC-048}: цель формулируется по шаблону
(анализировать \emph{объект} с целью \emph{назначение} относительно \emph{фокуса качества} с точки зрения
\emph{заинтересованного лица} в контексте \emph{контекста}), уточняется вопросами и только затем ---
метриками и данными. Метрика без цели не собирается. Определено восемь целей: G1 --- надёжность разметки и
согласие модели с учителями (RQ-1), G7 --- адаптация модели (RQ-2), G8 --- оценка усвоения (RQ-3), G2 --- как
учатся люди (RQ-4), а также описательные G3 --- объяснения (разведка), G4 --- вовлечённость и путь новичка,
G5 --- качество системы, G6 --- разработка с ИИ-агентами (раздел~\ref{sec:method-harness}).

Для пилота выбраны одна \emph{главная метрика} и одна \emph{контр-метрика}. Главная --- число строк книги,
впервые засчитанных учениками за неделю: она растёт, только когда ученики продвигаются по книге. Её можно
«накрутить», снизив порог, поэтому контр-метрика --- доля ложно засчитанных попыток: засчитанных
тренажёром, но не засчитанных по транскрипции учителя.

\section{Гипотезы}\label{sec:method-hyp}

Подтверждающие гипотезы перечислены в таблице~\ref{tab:hyp}. Заранее объявлены разведочными: сравнение
казахо- и русскоязычных на \ar{ق} и \ar{غ} (E-1), кривая «риск --- покрытие» (E-4), верификатор поверх
дообученной модели (E-5), нейросетевое отслеживание знаний (E-6), связь полного разбора ошибки с исходом
следующей попытки (E-7), модель забывания (E-8) и изменение доли ошибок чтения и произношения по книге (E-9).
Кривая обучения модели, вклад источников данных и кривые обучения учеников сообщаются обязательно, но
описательно. Всё записано в предрегистрации (раздел~\ref{sec:method-prov}); любой другой анализ помечается как
разведочный, а его выводы подаются как гипотезы для дальнейшей работы.

\begin{table}[ht]
\centering\small
\caption{Подтверждающие гипотезы (предрегистрация v0.4).}
\label{tab:hyp}
\begin{tabularx}{\linewidth}{@{}llL@{}}
\toprule
ID & RQ & Гипотеза \\
\midrule
H1.1 & RQ-1 & Учителя согласуются в типе ошибки (чтение / произношение) с $\kappa \ge 0{,}6$. \\
H1.2 & RQ-1 & $\kappa$ между B0 и учителем ниже $\kappa$ между учителями. \\
H1.3 & RQ-1 & У B0 ложные отклонения (FRR) встречаются чаще ложных принятий (FAR). \\
H1.4 & RQ-1 & Precision и recall обнаружения ошибок B0 на начинающих ниже, чем на домене обучения модели. \\
H2.1 & RQ-2 & На тестовых учениках FRR у B1 ниже, чем у B0, при recall не ниже recall B0 минус $\delta=0{,}05$. \\
H2.2 & RQ-2 & На тестовых учениках F1 по фонемам у B2 выше, чем у B0. \\
H2.3 & RQ-2 & На чтении экспертов PER у B2 выше, чем у B0, не более чем на 1 п.\,п. (не-худшесть). \\
H3.1 & RQ-3 & K2 поверх адаптированной модели предсказывает «золотой» исход следующей попытки с бо́льшим AUC,
чем K2 поверх B0. \\
H3.2 & RQ-3 & На одном распознавателе AUC у K2 выше, чем у K1. \\
H3.3 & RQ-3 & Итоговая оценка K2 согласуется со слепым тестом учителя: $\rho \ge 0{,}5$. \\
H4.1 & RQ-4 & Навыкам со звуками, которых нет в русском, до усвоения нужно больше попыток. \\
H4.2 & RQ-4 & У русскоязычных больше подтверждённых ошибок в твёрдости, чем в месте образования звука. \\
\bottomrule
\end{tabularx}
\end{table}

\section{Участники и материалы}\label{sec:method-participants}

Пилот закрытый: ученики приглашаются через преподавателей и знакомых автора, цель --- 20--40 взрослых
учеников в течение 4--6 недель. Участие в исследовании требует отдельного согласия; короткая анкета при
входе, которую видят только согласившиеся, фиксирует возраст (18+), родной язык и опыт чтения по-арабски.
Согласие на запись голоса явно разрешает обучение модели на записи и публикацию весов.
Не давшие согласия могут пользоваться тренажёром; их попытки оцениваются, но аудио не хранится, и в анализ
они не входят. Каждый ученик проходит короткий \emph{тест учителя} при входе, в конце пилота и через 1--2
недели после него: учитель слушает чтение фиксированного списка строк книги, покрывающего навыки
с.~10--24, и отмечает по каждому навыку «верно / неверно», не видя оценок системы. Строки теста не
используются для обучения моделей. Материал --- карточки уроков букв и огласовок (с.~10--24 книги), а для разведочного
анализа --- более поздние уроки, до которых дойдут ученики. Единица анализа --- попытка; ученики ---
кластеры, потому что попытки одного ученика зависимы.

Исключения зафиксированы заранее: попытки, оценённые заглушкой для разработки; записи, помеченные
разметчиком как непригодные или «читал не то» (считаются отдельно); попытки до заморозки
предрегистрации.

\section{Исследование 1: данные (RQ-1)}\label{sec:method-s1}

\paragraph{Разметка} Два преподавателя, уверенно читающих с таджвидом, размечают каждую попытку. Каждый
видит строку книги, запись, канонические фонемы, фонемы, услышанные тренажёром, и найденные им ошибки.
Если тренажёр услышал верно, разметчик ставит \emph{agree}; если нет --- \emph{fix} и вводит фонемы, которые
прозвучали на самом деле. Возможны ещё две метки: \emph{запись непригодна} и \emph{читал не то}. Отдельный
флаг отмечает \emph{самоисправление}, когда ученик ошибся и тут же перечитал верно в том же сообщении; по
таксономии Аль Мдфаа и соавторов \cite{SRC-068} такие попытки сообщаются отдельным блоком и не входят в
главные цифры RQ-1. Размечать можно в Telegram или в табличном пакете; оба пути пишут в одну таблицу с
кодом разметчика и версией руководства. Разметчики видят ответ тренажёра; это ускоряет работу, но может
склонять к согласию, что учитывается как угроза валидности.

\paragraph{Тип ошибки} Для каждой ошибки разметчик указывает тип. \emph{Ошибка чтения}: прочитан не тот знак ---
другая огласовка, пропущены сукун, шадда или мадд, перепутаны похожие на вид буквы (\ar{ب}, \ar{ت} и \ar{ث}).
\emph{Ошибка произношения}: знак прочитан верно, но звук произнесён неверно (место образования или признак).
Спорные случаи получают метку «не ясно». Оба учителя указывают тип на всех тестовых попытках и на каждой
десятой обучающей; надёжность --- $\kappa$ Коэна с «не ясно» как отдельной категорией и в анализе
чувствительности без неё (H1.1). Тип ошибки модели выводится правилом из выравнивания: замена огласовки, знака
или похожей на вид буквы --- чтение, замена по месту образования или признаку --- произношение; его согласие с
учителями сообщается описательно.

\paragraph{Метрики на уровне фонем} Для каждой попытки выравниваются каноническая, размеченная и
распознанная последовательности, и каждая каноническая фонема относится к TA, TR, FA или FR
(раздел~\ref{sec:lr-mdd}). Precision, recall, F1 и точность диагноза считаются по всем фонемам, по типу
ошибки и по букве \cite{SRC-064}.

\paragraph{Метрики на уровне вердикта} Вердикт учителя получают, применяя ту же политику оценки к его
транскрипции с порогами, сохранёнными в попытке. FAR --- доля попыток, которые тренажёр засчитал, хотя по
транскрипции учителя они не засчитываются; FRR --- доля не засчитанных, хотя по транскрипции они
засчитываются. Согласие измеряется $\kappa$ Коэна \cite{SRC-104}:
\begin{equation}
  \kappa = \frac{p_o - p_e}{1 - p_e},
\end{equation}
где $p_o$ --- наблюдаемое, а $p_e$ --- случайное согласие; $\kappa$ считается между тренажёром и каждым
учителем и между двумя учителями. Значение «учитель --- учитель» --- потолок, с которым сравнивается
тренажёр (H1.2). Согласие разметчиков по отдельным фонемам измеряется $\alpha$ Криппендорфа для
номинальных данных, допускающей пропуски и любое число оценщиков \cite{SRC-065}.

\paragraph{Неопределённость} 95-процентные доверительные интервалы строятся кластерным бутстрэпом с
выборкой учеников с возвращением. Сравнение с опубликованными бенчмарками --- только качественное, потому
что данные разные.

\paragraph{Профиль ошибок} Частоты типов ошибок, матрицы путаниц «ожидалось $\to$ услышано» и трудность
по строкам книги и по её минимальным парам считаются дважды: по всем ошибкам, найденным тренажёром, и по
ошибкам, подтверждённым разметчиками (верные отклонения). Разница показывает, какую часть видимого профиля
ученика создаёт модель. Чтобы выводы не определялись одним активным учеником, каждая частота сообщается
вместе с числом учеников, которые в неё вносят вклад. Подтверждённый профиль ошибок --- данные для
исследования~4.

\section{Исследование 2: как учится ИИ (RQ-2)}\label{sec:method-s2}

\paragraph{Системы} Сравниваются три системы на одних и тех же записях. \textbf{B0} --- Muaalem как есть с
выравниванием и порогами (глава~\ref{ch:impl}). \textbf{B1} --- верификатор поверх выходов B0. Для каждой
канонической фонемы $j$ попытки он получает вектор признаков $x_j$: уверенность модели по фонеме,
вероятности артикуляционных признаков, меру Goodness of Pronunciation
\begin{equation}
  \mathrm{GOP}(p_j) = \frac{1}{d_j}\log\frac{P(p_j\mid O_j)}{\max_{q} P(q\mid O_j)},
\end{equation}
где $O_j$ --- кадры звука, отнесённые к фонеме, $d_j$ --- их число \cite{SRC-115}, а также длительность, тип
ошибки по выравниванию и тип материала (Коран или нет). Логистическая модель или градиентный бустинг
оценивает $\hat p_j = P(\text{ошибка}\mid x_j)$; вероятности калибруются на отдельных учениках обучающей
части, качество калибровки измеряется ожидаемой ошибкой калибровки (ECE) \cite{SRC-112}. Попытка
отклоняется, если $\max_j \hat p_j \ge \lambda$. Порог $\lambda$ выбирается конформным контролем риска
\cite{SRC-114}: на калибровочных учениках берётся наименьший $\lambda$, при котором
\begin{equation}
  \frac{n}{n+1}\,\widehat{\mathrm{FRR}}_n(\lambda) + \frac{1}{n+1} \le \varepsilon,
\end{equation}
где $\widehat{\mathrm{FRR}}_n$ --- доля ложных отклонений на $n$ калибровочных попытках, $\varepsilon=0{,}05$
(QAS-02). Попытки с $\hat p_j$ в зоне неопределённости вокруг порога получают исход «не уверен»
\cite{SRC-117}: ученику предлагают перечитать, попытка уходит учителю. Гарантия конформного метода
требует перестановочности; попытки одного ученика зависимы, поэтому калибровка ведётся по ученикам, а
выполнение гарантии на тесте проверяется, а не предполагается.

\textbf{B2} --- Muaalem, дообученная методом LoRA \cite{SRC-111}: к матрицам внимания $W_0\in\mathbb{R}^{d\times k}$
добавляется поправка низкого ранга
\begin{equation}
  W = W_0 + \tfrac{\alpha}{r}\,BA, \qquad B\in\mathbb{R}^{d\times r},\; A\in\mathbb{R}^{r\times k},\; r\ll\min(d,k),
\end{equation}
обучаются только $A$, $B$ и выходные CTC-головы, $W_0$ заморожены. Цель обучения --- многоуровневый CTC
Muaalem на фонемах, исправленных преподавателями. Обучающие данные пополняются двумя способами: по
минимальным парам книги (верное чтение карточки с одной буквой пары, подписанное каноническими фонемами
парной карточки, --- пример замены) \cite{SRC-119} и изменением скорости $0{,}9/1{,}0/1{,}1$
\cite{SRC-120}. Сетка гиперпараметров ($r\in\{8,16\}$, скорость обучения, число эпох с ранней остановкой)
фиксирована в предрегистрации.

\paragraph{Разбиение и защита теста} До первого обучения около 30\,\% учеников, стратифицированно по
родному языку и активности, уходят в тест; список фиксируется отпечатком SHA-256. Разбиение --- по
ученикам, а не по попыткам: иначе голос одного диктора оказался бы и в обучении, и в проверке
\cite{SRC-122}. Гиперпараметры B2 и калибровка B1 выбираются только 5-кратной кросс-валидацией по
ученикам внутри обучающей части \cite{SRC-121}. Тестовые попытки оба преподавателя размечают вслепую,
не видя ответа тренажёра, и согласуют расхождения; это «золото», с которым сравниваются все три системы.
Тест открывается один раз, когда все системы обучены и заморожены, а оценщик зафиксирован отпечатком.

\paragraph{Метрики и анализ} Для каждой системы на тесте считаются метрики исследования~1: precision,
recall, F1 и точность диагноза по фонемам, FAR и FRR по вердикту, $\kappa$ с «золотом». H2.1 и H2.2
проверяются парным кластерным бутстрэпом по ученикам: разница метрик двух систем пересчитывается на
одних и тех же выборках учеников. H2.3 проверяется на отложенном наборе чтения экспертов, не входившем в
обучение: верхняя граница 95-процентного интервала разницы PER должна быть меньше одного процентного
пункта \cite{SRC-118}. Сводная мера --- \emph{доля закрытого разрыва}
\begin{equation}
  G_B = \frac{\kappa_B - \kappa_{B0}}{\kappa_{TT} - \kappa_{B0}},
\end{equation}
где $\kappa_{TT}$ --- согласие учителей между собой: $G_B=0$ --- адаптация ничего не дала, $G_B=1$ ---
система согласуется с учителем так же, как учителя друг с другом. Обязательно сообщаются \emph{кривая
обучения} B1 и B2 (25, 50, 75 и 100\,\% обучающих учеников), \emph{вклад источников данных} --- только поправки
учителей; плюс минимальные пары книги; плюс оспаривания учеников «бот ошибся» --- и доля попыток, которые
система решает сама при риске не выше $\varepsilon$, то есть сколько работы она снимает с учителя. Разведочно
--- кривая «риск --- покрытие» B1.

\paragraph{Запасной путь} Если к 15 апреля 2027~г. в обучающей части меньше примерно 1\,500 размеченных
попыток, B2 подаётся как проверка осуществимости, H2.2 и H2.3 снимаются новой версией предрегистрации
до открытия теста, а RQ-2 отвечает B1. Решение принимается по числу размеченных попыток, а не по
результатам.

\paragraph{Теневой запуск} После обучения B1 и B2 работают в тренажёре в тени: ученик видит вердикт B0,
вердикты адаптированных систем записываются в попытку. Ученики весь пилот видят одну и ту же систему, поэтому
смена модели не смешивается с их продвижением (RQ-4), а попытки после обучения дают описательную проверку.

\section{Исследование 3: измерение усвоения (RQ-3)}\label{sec:method-s3}

\paragraph{Карта навыков} Навык --- соответствие «буква + огласовка или знак $\to$ звук», которое вводит книга;
навыки упорядочены так, как их вводят уроки. Каждая строка книги проверяет несколько навыков; матрица
$Q=(q_{jk})$, где $q_{jk}=1$, если строка $j$ проверяет навык $k$, строится из программы автоматически
(глава~\ref{ch:impl}). Исход попытки по навыку --- верно ли прочитаны слоги этого навыка.

\paragraph{Модели} \textbf{K0} --- доля верных попыток ученика по навыку. \textbf{K1} --- байесовское
отслеживание знаний \cite{SRC-124}: навык ученика --- скрытое состояние $L_t\in\{0,1\}$ («усвоен» или нет),
наблюдение $o_t$ --- верен ли исход попытки $t$, параметры --- начальная вероятность $P(L_0)$, вероятность
научиться $T$, угадывания $g$ и оплошности $s$:
\begin{equation}
  P(o_t{=}1\mid L_t{=}1)=1-s,\qquad P(o_t{=}1\mid L_t{=}0)=g,
\end{equation}
\begin{equation}
  P(L_{t+1}{=}1) = P(L_t{=}1\mid o_t) + \bigl(1-P(L_t{=}1\mid o_t)\bigr)\,T .
\end{equation}
\textbf{K2} учитывает, что исход наблюдает не учитель, а распознаватель, который сам ошибается: наблюдается
вердикт системы $\hat o_t$, а не истинный исход $o_t$. Для фонемы $p$ навыка ложные принятия
$\mathrm{FAR}_p$ и ложные отклонения $\mathrm{FRR}_p$ распознавателя измеряются против учителей на обучающей
части (исследования~1 и~2), и
\begin{equation}
  P(\hat o_t{=}1\mid L_t) = P(o_t{=}1\mid L_t)\,(1-\mathrm{FRR}_p) + P(o_t{=}0\mid L_t)\,\mathrm{FAR}_p .
\end{equation}
Так шум распознавателя не растворяется в $g$ и $s$, как у Бека и Сисона \cite{SRC-125}, а входит в модель как
измеренная величина, и адаптация распознавателя (меньшие $\mathrm{FAR}_p$, $\mathrm{FRR}_p$) должна делать
оценку усвоения точнее. Параметры оцениваются на обучающих учениках (EM или поиск по сетке); K1 и K2 строятся
поверх B0 и поверх адаптированной системы, выбранной кросс-валидацией на обучении.

\paragraph{Проверка} На тестовых учениках модели предсказывают «золотой» исход следующей попытки по навыку
(разметка учителей вслепую, а не вердикт системы). Метрики --- AUC, логарифмическая потеря и ошибка калибровки
\cite{SRC-126}. H3.1 и H3.2 проверяются парным кластерным бутстрэпом разницы AUC по ученикам. Внешний критерий
(H3.3) --- тест учителя в конце пилота: итоговая оценка усвоения K2 по группам навыков сравнивается с долей
верных навыков группы по тесту $\rho$ Спирмена, интервал --- бутстрэпом по ученикам. Разведочно сравнивается
нейросетевое отслеживание знаний \cite{SRC-127}; на малых данных от него не ждут выигрыша.

\section{Исследование 4: как учатся люди (RQ-4)}\label{sec:method-s4}

\paragraph{Кривые обучения} Для каждого навыка строится доля ошибок в зависимости от номера попытки ученика по
этому навыку. Модель аддитивных факторов \cite{SRC-128}
\begin{equation}
  \operatorname{logit} P(Y_{ij}{=}1) = \theta_i + \sum_k q_{jk}\,(\beta_k + \gamma_k N_{ik}),
\end{equation}
где $\theta_i$ --- способность ученика $i$, $\beta_k$ --- лёгкость навыка $k$, $\gamma_k$ --- скорость его
освоения, $N_{ik}$ --- число предыдущих попыток ученика по навыку, показывает, какие навыки даются быстро
(большое $\gamma_k$) и какие медленно. Она же проверяет карту навыков: если навык выбран неудачно, его кривая не
падает. Подгонка AFM с картой книги сравнивается с картой «одна буква --- один навык».

\paragraph{Попытки до усвоения и забывание} Навык считается усвоенным, когда оценка K2 достигает
$P(L)\ge0{,}95$. Число попыток до усвоения сравнивается для навыков со звуками, которых нет в русском, и
остальных (H4.1) смешанной моделью со случайным эффектом ученика \cite{SRC-106}. H4.2 использует выходы модели
по признакам (\emph{сифат}) на подтверждённых ошибках и проверяется смешанной логистической моделью. Забывание
описывается двумя способами: удержание на отложенном тесте учителя и падение вероятности верного исхода после
перерыва между занятиями; модель периода полураспада памяти \cite{SRC-129} --- разведка.

\paragraph{Тест учителя} Доля верных навыков при входе, в конце пилота и на отложенном тесте сообщается
описательно, по группам навыков. Без контрольной группы рост между тестами не доказывает, что учеников научила
система: это описание того, как учатся люди с системой. По образцу DART \cite{SRC-130,SRC-131} точность ИИ и
обучение людей сообщаются отдельно. По двенадцати подтверждающим гипотезам процедура Холма контролирует
групповую ошибку на уровне $\alpha=0{,}05$; размеры эффекта и доверительные интервалы сообщаются всегда. Сбор
данных останавливается в запланированный срок окончания пилота; промежуточные отчёты --- только описательные.

\section{Достоверность разработки и экспериментов}\label{sec:method-harness}

Систему создаёт один автор вместе с ИИ-агентом, поэтому её числа должны быть защищены от подгонки не только
со стороны автора, но и со стороны агента. Обвязка разработки (раздел~\ref{sec:impl-harness}) описывается
как средство достоверности и документируется по цели G6 как встроенное кейс-стади \cite{SRC-102}: случай ---
разработка системы; единицы анализа --- рабочие сессии и отдельные изменения. Источников данных четыре: журнал
хуков, фиксирующий каждую блокировку и запрос подтверждения со временем, хуком и причиной; сводки сессий;
журнал каждого запроса к агенту; история версий с ADR. Метрики: блокировки и запросы подтверждения по
причинам; вмешательства человека за сессию; итерации до прохождения проверок; дефекты, прошедшие все проверки;
попытки агента обратиться к тестовой выборке. Самооценка продуктивности не используется \cite{SRC-030}. Станет
ли эта часть отдельным исследовательским вопросом или отдельной статьёй, решается с руководителем.

\section{Предрегистрация, происхождение данных и воспроизводимость}\label{sec:method-prov}

Три механизма защищают результаты от подгонки задним числом.

\paragraph{Предрегистрация} Гипотезы, исключения, объём выборки и анализ записываются до сбора данных.
Когда автор и руководитель утверждают план, фиксируются его отпечаток SHA-256 и дата; любое последующее
изменение требует новой версии с указанием причины. Отдельно, до первого обучения, отпечатками фиксируются
список тестовых учеников, код оценщика RQ-2 и RQ-3 и список строк теста учителя. Заморозка данных и оценщика хешем следует практике группы
научного руководителя \cite{SRC-067}.

\paragraph{Защищённый оценщик и тест} Модуль, считающий метрики RQ-1--RQ-3, --- зона человека: агент должен
спросить, прежде чем менять его, потому что агент, способный править оценщик, может улучшить счёт, не
улучшив систему \cite{SRC-027,SRC-082}. Метки тестовых учеников агенту читать запрещает хук: агент,
видевший тест, мог бы подобрать под него модель.

\paragraph{Происхождение} В каждой попытке хранятся канонические варианты, распознанные фонемы, сырой
выход модели, пороги, версии кода, программы и модели, строка книги в том виде, в каком её видел ученик, и
SHA-256 аудио; у дообученной модели --- хеш весов и хеш обучающей выборки. Команда трассировки печатает отпечатки кода, программы, SQL-отчётов, миграций, зависимостей
и предрегистрации. Число попадает в диссертацию, только если соответствующий анализ реализован, а у рисунка
или таблицы есть паспорт: отчёт, дата, число попыток и учеников, отпечаток кода.

\section{Этика}\label{sec:method-ethics}

Исследование связано с голосовыми записями на религиозную тему, поэтому оно спроектировано так, будто
действует самое строгое прочтение закона. Согласие на работу сервиса и согласие на запись голоса
разделены; аудио без согласия не хранится; согласие можно отозвать, что удаляет записи и разметку по ним;
каждый ученик может удалить все данные одной кнопкой; в отчётах и наборах данных --- только псевдонимы;
участие в исследовании --- только для взрослых; голосовые отпечатки и эмбеддинги диктора не вычисляются;
согласие на запись голоса явно разрешает обучение модели и публикацию весов, а отзыв согласия исключает
записи из следующих версий набора и модели;
разметчики видят только строку и ответ тренажёра, но не личность ученика, а записи им отправляются в
защищённом режиме \cite{SRC-033,SRC-036,SRC-037}. Вопросы о том, кто оператор данных, нужна ли
экспертиза этической комиссии и где расположен сервер, --- институциональные и решаются с руководителем до
приглашения учеников вне закрытого круга.

\section{Угрозы валидности}\label{sec:method-validity}

Таблица~\ref{tab:validity} перечисляет главные угрозы по классификации Волина и соавторов
\cite{SRC-106} и то, как на них отвечает дизайн.

\begin{table}[ht]
\centering\small
\caption{Угрозы валидности и ответ на них.}
\label{tab:validity}
\begin{tabularx}{\linewidth}{@{}P{2.4cm}LL@{}}
\toprule
Вид & Угроза & Ответ \\
\midrule
Конструктная & Метка разметчика --- не «истинная» ошибка; учителя могут расходиться. & Два разметчика;
$\kappa$ «учитель --- учитель» как потолок; письменное руководство с версиями. \\
Конструктная & Разметчики видят ответ тренажёра и могут соглашаться по подсказке, что тянет сравнение к
B0. & Тест размечается вслепую; H1.3 сообщается и на слепой части; согласие анализируется отдельно
для засчитанных и не засчитанных попыток. \\
Внутренняя & Утечка между обучением и тестом (один диктор, подбор под тест). & Разбиение по ученикам;
список теста заморожен; CV внутри обучения; хук на метки теста; тест открывается один раз. \\
Внутренняя & Выигрыш на новичках куплен забыванием. & H2.3 на отложенном чтении экспертов. \\
Выводов & Тест мал (около десяти учеников). & Парный кластерный бутстрэп; широкие интервалы
сообщаются как есть; кривая обучения показывает, сколько данных не хватает. \\
Конструктная & Граница «ошибка чтения / произношения» зависит от правила разметки. & H1.1 проверяет её
надёжность; «не ясно» --- отдельная категория; правило с примерами по книге. \\
Конструктная & Навыки книги могут быть выбраны неудачно. & Кривые AFM по навыкам; сравнение с картой «одна
буква --- один навык». \\
Внутренняя & Прогресс учеников может объясняться не системой (занятия с учителем, повтор). & RQ-4
формулируется как описание; тест учителя вслепую; без причинных выводов. \\
Внутренняя & Тест учителя проводит тот же учитель, что размечает. & Тест вслепую к оценкам системы; список строк
фиксирован заранее и не входит в обучение. \\
Внутренняя & Длинные строки «разбавляют» ошибки в мере сходства. & RQ-1 считается по фонемам, а не только
по вердикту. \\
Внутренняя & Агент может подогнать собственные проверки. & Замороженные тесты, защищённый оценщик, журнал
блокировок (раздел~\ref{sec:method-harness}). \\
Внешняя & Закрытый пилот из знакомых; малая выборка. & Подаётся как пилот; размеры эффекта с
интервалами; без обобщения за пределы аудитории. \\
Внешняя & Модель обучена на Коране, а с.~10--37 не из Корана. & Результаты --- по типу материала. \\
Выводов & Зависимые попытки; множественные проверки. & Кластерный бутстрэп; смешанные модели; поправка
Холма. \\
\bottomrule
\end{tabularx}
\end{table}
